\documentclass[12pt]{article}
\usepackage[utf8]{inputenc}
\usepackage[T1]{fontenc}
\usepackage{hyperref}
\title{Tunisian Sign Language Language Resources for K-12 Curriculum Development}
\author{Compiled by Claude Sonnet 5 (claude-sonnet-5) for Andr\'as Kornai}
\date{}
\begin{document}
\maketitle

\section{Introduction}
Tunisian Sign Language (LST, ISO 639-3 tse) is the sign language of Tunisia's deaf community. Deaf education there is still predominantly oralist. A 2023 Frontiers article reports that no formal LST dictionary exists and that the lack of suitable LST teaching materials led many teachers in deaf institutions to use French Sign Language (LSF) instead. Search results were mostly about SignWriting work and academic papers.

\section{Primary Educational Resources}
\begin{itemize}
\item Balti, Bahri and Laajili describe using SignWriting in mainstream Tunisian schools. They translated some primary-school modules (for example mathematics and science) into LST and taught deaf children with it: \url{https://signwriting.org/symposium/archive/sws0015_Slides_Tunisia_SignWriting_Deaf_Education_Balti_Bahri_Laajili.pdf}. The abstract is at \url{https://www.signwriting.org/symposium/archive/sws0015_Abstract_Tunisia_SignWriting_Deaf_Education_Balti_Bahri_Laajili.pdf}.
\item The SignWriting Tunisia page lists a free Tunisian SignWriting workbook (French and Arabic) and Goldilocks storybooks in LST at levels 1--3: \url{https://signwriting.org/tunisia}. I did not open the individual files.
\end{itemize}

\section{Linguistic and Reference Materials}
\begin{itemize}
\item The SignWriting Tunisia page links an online LST dictionary on SignPuddle, started by Mohamed Ali Balti. It is described as a continuing project: \url{https://signwriting.org/tunisia}. I did not check whether it is currently live.
\item Deafness in Tunisia (overview): \url{https://en.wikipedia.org/wiki/Deafness_in_Tunisia}
\item Frontiers in Education, 2021, a proposed bicultural trilingual education model for deaf children in Tunisia: \url{https://www.frontiersin.org/articles/10.3389/feduc.2021.750584/pdf}
\item Frontiers in Communication, 2023, the article reporting the dictionary and materials gap: \url{https://www.frontiersin.org/journals/communication/articles/10.3389/fcomm.2023.1201148/pdf}
\end{itemize}

\section{Religious and Community Resources}
AVST (Association Voix du Sourd de Tunisie) is described in the Wikipedia article as the oldest organisation for the deaf in Tunisia. See \url{https://en.wikipedia.org/wiki/Deafness_in_Tunisia}. No website was found, and no contact details are given here.

\section{Digital and Audio Resources}
The SignPuddle dictionary and the SignWriting storybooks (above) are the only digital resources found. No video corpus was located.

\section{Curriculum Development Context}
A practical base could combine the SignWriting-based translated modules, the workbook, the dictionary and the trilingual model proposal. Pedagogically adequate LST materials remain scarce, and the existing materials are small in scale.
\end{document}